\chapter{CONCLUSION}

\section{Introduction}
This chapter concludes the research by summarizing the key findings, drawing 
empirical conclusions, outlining the limitations of the current work, and 
identifying potential directions for future research. The core focus of this 
dissertation was to address the inefficiencies of standard dynamic array 
resizing policies, which rely on rigid, workload-independent growth factors 
(e.g., $1.5\times$ or $2.0\times$). By introducing runtime feedback loops that 
incorporate resize history ($k$) and real-time insertion velocity ($v$), this 
study successfully developed and optimized an adaptive resizing framework that 
balances memory efficiency and computational speed.

\section{Summary of Findings}
The empirical results compiled across $10^7$ element benchmarks (detailed in 
Chapter 4) demonstrate the high efficiency of the proposed workload-aware 
resizing models. The key findings of this research are summarized below:
\begin{itemize}
    \item \textbf{Optimal Overall Model:} The GP-Optimized Hybrid Framework 
    (Formula D) achieved the lowest overall cost score, $J = 0.2589$. With optimized 
    parameters ($w_A = 0.9955$, $\lambda = 3.0944$), Formula D restricted 
    memory waste to a mere \textbf{0.42\%} while requiring only \textbf{20 resizes} 
    and \textbf{10.49M elements copied}. This matched the optimal copy speed 
    of the C++ STL doubling policy ($2.0\times$) while reducing final memory 
    waste by \textbf{90.9\%} (from 4.63\% down to 0.42\%).
    \item \textbf{Comparison with Static Standards:} Compared to other industry 
    standards, the proposed Formula D demonstrated substantial improvements:
    \begin{itemize}
        \item It reduced memory waste by \textbf{98.5\%} compared to Java's 
        $1.5\times$ policy (27.77\% waste) and cut element copying overhead by 
        \textbf{62.1\%} (from 27.69M copies down to 10.49M).
        \item It reduced cumulative element copies by \textbf{87.5\%} compared 
        to CPython's list policy (83.84M copies down to 10.49M) while maintaining 
        comparable memory utilization.
    \end{itemize}
    \item \textbf{Absolute Memory Waste Leader:} The Bayesian-tuned Formula A 
    Phase 1 (Logarithmic) achieved the lowest absolute memory waste of \textbf{0.05\%} 
    at the end of insertion, demonstrating the theoretical limit of capacity-based 
    decay, albeit at the cost of 32.31M element copies.
    \item \textbf{Workload Adaptability:} The feedback-enabled models successfully 
    adapted to diverse data ingestion profiles. By incorporating velocity $v$, the 
    array expanded aggressively during high-speed burst phases to skip costly 
    reallocations, and grew conservatively during slow, sparse insertion phases 
    to prevent over-allocation.
\end{itemize}

\section{Conclusion Drawn by Results}
Based on the empirical evidence, several conclusions can be drawn:
\begin{enumerate}
    \item \textbf{Effectiveness of the Proposed Strategy:} The research successfully 
    proved that incorporating runtime feedback of resize history 
    and insertion velocity allows dynamic arrays to achieve a superior trade-off 
    between spatial and temporal efficiency compared to static, fixed growth 
    factors.
    \item \textbf{Harmonic Decay is Highly Effective:} Decaying the growth 
    factor harmonically based on the resize counter ($k$) is an extremely effective 
    mechanism for contiguous arrays. It naturally limits the absolute volume of 
    wasted space as the array grows massive, preventing the geometric waste 
    accumulation seen in standard doubling policies.
    \item \textbf{Throughput Preservation:} The computational overhead of evaluating 
    logarithmic and harmonic math operations during resize events is negligible 
    when compared to the massive time savings achieved by avoiding memory allocation 
    calls and element copying loops. Thus, adaptive resizing improves overall 
    efficiency without sacrificing throughput.
\end{enumerate}

\section{Limitations}
While the proposed algorithm demonstrated high efficiency, several limitations 
were identified during this study:
\begin{itemize}
    \item \textbf{Timestamp Measurement Overhead:} Querying the high-resolution 
    system clock (\texttt{time.perf\_counter()}) to compute insertion velocity 
    introduces a small but non-zero CPU overhead during each resize cycle. In 
    extremely high-performance systems with tiny arrays, this system call latency 
    can slightly degrade throughput.
    \item \textbf{Simulation Boundaries:} The evaluation was conducted inside 
    a custom Python-based simulation environment. Although the simulator accurately 
    models logical capacity, size, and element copying overheads, it operates 
    within CPython's managed memory heap. The actual physical memory layout and 
    virtual memory fragmentation are thus dependent on CPython's garbage collector, 
    introducing minor measurement noise.
    \item \textbf{Static Clamping Bounds:} The lower and upper clamping bounds 
    ($g_{\min} = 1.1$, $g_{\max} = 2.0$) were set statically. In systems with 
    highly unusual memory pages, these bounds may need to be adjusted dynamically.
\end{itemize}

\section{Future Works}
To address the limitations and expand the applicability of this research, 
the following areas are proposed for future study:
\begin{itemize}
    \item \textbf{Native C++ STL Integration:} Porting the adaptive growth 
    engine directly into a custom C++ STL allocator or modifying the 
    \texttt{std::vector} implementation. Testing in a compiled language will 
    eliminate runtime interpreter noise and allow measuring raw, bare-metal CPU 
    execution cycles.
    \item \textbf{Deployment in Embedded Systems (IoT):} Evaluating the 
    algorithm on resource-constrained microcontrollers (e.g., ESP32, ARM Cortex-M) 
    where memory is highly restricted (measured in kilobytes) and memory 
    fragmentation can lead to system crashes.
    \item \textbf{Predictive Machine Learning Models:} Replacing the simple 
    historical velocity averages with lightweight online regression or reinforcement 
    learning models that can predict upcoming data bursts, allowing the array to 
    pre-allocate memory preemptively.
    \item \textbf{Multi-Threaded Vector Optimizations:} Extending the velocity 
    computations to support thread-safe, lock-free concurrent vector structures, 
    adjusting the growth factor based on collective thread insertion speeds.
\end{itemize}

\section{Summary for the Chapter}
This chapter concluded the research by summarizing the findings and drawing 
conclusions on workload-aware dynamic array resizing. Section 5.2 summarized the 
empirical findings, highlighting that the GP-Optimized Hybrid Framework 
(Formula D) achieved a Pareto-optimal J-Score of 0.2589, cutting memory waste 
by 90.9\% compared to C++ doubling while matching its copy speed. Section 5.3 
validated the core design strategy, confirming that harmonic resize decay 
coupled with velocity scaling balances space and time effectively. Section 5.4 
discussed current limitations, including clock measurement overhead and simulation 
boundaries. Finally, Section 5.5 outlined future directions, including C++ STL 
vector integration, embedded IoT deployments, and predictive models, providing 
a clear roadmap for subsequent research.
